\subsection{\textit{Computer Vision} dan Citra Digital}

\textit{Computer Vision} digunakan untuk memperoleh informasi dari citra melalui proses komputasi. Pada analisis makanan, informasi visual dapat dimanfaatkan untuk melakukan pengenalan, deteksi, dan segmentasi objek sehingga jenis serta lokasi makanan pada citra dapat ditentukan \cite{aguilar2017}.

Citra digital berwarna tersusun atas nilai piksel pada beberapa kanal warna. Informasi tersebut menjadi masukan bagi model untuk membentuk representasi fitur yang berkaitan dengan karakteristik visual objek, seperti bentuk, warna, tekstur, dan susunan objek pada citra. Pada sistem yang dirancang, citra piring digunakan sebagai masukan bagi YOLO11-seg untuk mengenali lauk pada tingkat instans.

Karena pengambilan citra pada sistem dilakukan dari posisi tetap di atas piring, area objek dapat diamati dengan perspektif yang relatif konsisten. Pendekatan pengambilan citra dari arah atas juga digunakan pada sistem \textit{visual checkout} makanan untuk memperoleh tampilan objek pada area pengamatan secara konsisten \cite{wu2021}.
